\SourceLeaf{303}%
Punkte $A$, $B$, $C$, $D$ (ich verwende hier meine eigene Bezeichnung, da ich \textsc{Newtons}
in diesem Fall für unhandlich halte) einen Punkt $\delta$ jener Geraden
sucht und dann mithilfe von $A$, $B$, $C$, $E$ einen anderen Punkt $\varepsilon$ jener
Geraden. Damit ist die Gerade selbst der Lage nach gegeben. Für jeden anderen
Punkt auf ihr, $\mu$, liefert seine Konstruktion einen entsprechenden
Punkt $M$ auf dem Kegelschnitt. Er lässt den Winkel $BAC$ (mit unbestimmter Schenkellänge)
um $A$ rotieren und hält die Stellung fest, in der $AB$ durch $\mu$ geht. Dann muss $AC$
durch $M$ gehen. Dreht sich ebenso der Winkel $ABC$ um $B$, so muss, wenn
$BA$ durch $\mu$ geht, auch $BC$ durch $M$ gehen. Diese Konstruktion ist
überaus elegant, aber Herr der Sache ist man erst, wenn man
alle anderen Fragen zu dieser magischen Geraden beantworten kann.
Namentlich will man wissen, in welchen Beziehungen diese Gerade zu den Elementen des
Kegelschnitts steht und ob sich diese Elemente selbst mit Leichtigkeit aus der Lage jener
Geraden und der Lage von $A$, $B$, $C$ ableiten lassen. Manches dieser Art kann
ich jetzt recht elegant ausführen. Ich weiß aber nicht, ob ich das Ganze selbst
durchführen kann, da mich andere Geschäfte zwingen, abzubrechen.

\mbox{.\,.\,.\,.\,.\,.\,.} Früher haben Sie sich, wenn ich mich nicht irre, viel mit den Kegelschnitten
beschäftigt und gewiss viel mehr darüber gelesen als ich. Sie führen zum Beispiel schon in
Monatliche Correspondenz, 1810 November, S. 508, einen eleganten Satz über die Ellipse
aus einem Buch\HistoricalNote{[*)]}{[*)}{\textsc{Puissant}, Recueil de diverses propositions de géométrie. Paris 1801.
Vgl. auch Bd. IV, S. 396.]} an, das ich bisher noch nicht angesehen habe, obwohl unsere Bibliothek es
vermutlich besitzt. Wegen historischer Hinweise kann man also wohl zuerst bei Ihnen nachfragen.
Kennen Sie eine elegante Konstruktion, um den \textls{Mittelpunkt}
eines Kegelschnitts aus 5 seiner Punkte --- \textls{ohne alle überflüssigen Umwege} --- zu
finden? Schwer wird es gewiss nicht sein, aber gerade diese Herleitung habe ich noch
nicht \mbox{ausgeführt. .\,.\,.\,.\,.\,.\,.}
